\chapter{Learning to Rank: Fondamenti, Tassonomia e Ottimizzazione di BM25F}
\label{chap:learning-to-rank-fondamenti}

\FloatBarrier
\section{Introduzione: Limiti del Bag of Words e Framework Generale di Ranking}
\label{sec:ltr-intro-framework-generale}

I modelli di reperimento tradizionali basati sul paradigma Bag of Words (BoW), quali il Vector Space Model e Okapi BM25, hanno costituito il pilastro fondante dei motori di ricerca commerciali. Tuttavia, la crescente disponibilit\`a di segnali eterogenei e la necessit\`a di combinare evidenze lessicali, topologiche e comportamentali hanno reso indispensabile l'adozione di metodologie fondate sul Machine Learning.

Questa sezione formalizza la decomposizione astratta delle funzioni di ranking basata sulle rappresentazioni~\cite{guo2020}, analizza l'insufficienza dei vettori sparsi BoW e illustra la ricchezza dello spazio dei segnali moderni.

\subsection{Decomposizione Funzionale delle Funzioni di Ranking}
\label{subsec:decomposizione-funzionale-ranking}

Seguendo la formalizzazione proposta da Guo et al.~\cite{guo2020}, una funzione di ranking generale che assegna un punteggio di rilevanza $s(q, d) \in \mathbb{R}$ a una coppia query-documento $(q, d)$ pu\`o essere decomposta in tre moduli di rappresentazione e un operatore di aggregazione finale.

\begin{definition}[Framework Generale di Rappresentazione per il Ranking]
Siano dati una query $q \in \mathcal{Q}$ emessa da un utente e un documento $d \in \mathcal{D}$ appartenente alla collezione indicizzata. La computazione dello score di rilevanza $s(q, d)$ \`e descritta dalla composizione:
\begin{equation}
    s(q, d) = f\big(\phi(q), \psi(d), \eta(q, d)\big),
    \label{eq:general-ranking-framework}
\end{equation}
dove:
\begin{itemize}
    \item $\phi: \mathcal{Q} \to \mathcal{X}_q$ definisce la funzione di \textbf{rappresentazione della sola query};
    \item $\psi: \mathcal{D} \to \mathcal{X}_d$ definisce la funzione di \textbf{rappresentazione del solo documento};
    \item $\eta: \mathcal{Q} \times \mathcal{D} \to \mathcal{X}_{qd}$ definisce la funzione di \textbf{interazione query-documento};
    \item $f: \mathcal{X}_q \times \mathcal{X}_d \times \mathcal{X}_{qd} \to \mathbb{R}$ rappresenta la \textbf{funzione di aggregazione} che sintetizza i vettori nello scalare finale di ordinamento.
\end{itemize}
\end{definition}

\begin{figure}[htbp]
\centering
\begin{tikzpicture}[
    box/.style={
        rectangle, 
        draw=navy, 
        fill=codebackground, 
        rounded corners=4pt, 
        line width=0.8pt, 
        align=center, 
        text width=3.4cm, 
        minimum height=1.1cm,
        inner sep=5pt, 
        font=\small
    },
    inputbox/.style={
        rectangle, 
        draw=navy, 
        fill=codebackground, 
        rounded corners=4pt, 
        line width=0.8pt, 
        align=center, 
        text width=2.8cm, 
        minimum height=0.9cm,
        inner sep=4pt, 
        font=\small\bfseries
    },
    opbox/.style={
        rectangle, 
        draw=darkgreen, 
        fill=darkgreen!10, 
        rounded corners=4pt, 
        line width=0.9pt, 
        align=center, 
        text width=5.8cm, 
        minimum height=1.2cm,
        inner sep=6pt, 
        font=\small\bfseries
    },
    scorebox/.style={
        rectangle, 
        draw=crimson, 
        fill=crimson!10, 
        rounded corners=4pt, 
        line width=0.9pt, 
        align=center, 
        text width=4.6cm, 
        minimum height=1.1cm,
        inner sep=6pt, 
        font=\small\bfseries
    },
    arrow/.style={->, >=stealth, line width=0.9pt, navy}
]
    % Riga 1: Ingressi (distanza orizzontale 8.4 cm)
    \node[inputbox] (q) at (0, 4.0) {Query $q$};
    \node[inputbox] (d) at (8.4, 4.0) {Documento $d$};

    % Riga 2: Rappresentazioni intermedie (spaziate a 0 cm, 4.2 cm, 8.4 cm)
    \node[box] (phi) at (0, 2.0) {Rappresentazione Query\\[2pt]$\phi(q)$};
    \node[box] (eta) at (4.2, 2.0) {Rappresentazione\\Query--Documento\\[2pt]$\eta(q, d)$};
    \node[box] (psi) at (8.4, 2.0) {Rappresentazione\\Documento\\[2pt]$\psi(d)$};

    % Riga 3: Funzione di Aggregazione (centrata sotto a 4.2 cm, y = -0.4)
    \node[opbox] (f) at (4.2, -0.4) {Funzione di Aggregazione\\[3pt]$f\big(\phi(q), \psi(d), \eta(q, d)\big)$};

    % Riga 4: Score di Rilevanza finale (y = -2.4, spaziatura verticale generosa di 2.0 cm)
    \node[scorebox] (score) at (4.2, -2.3) {Punteggio di Rilevanza Finale\\[3pt]$s(q, d)$};

    % Connessioni dagli ingressi alle rappresentazioni
    \draw[arrow] (q) -- (phi);
    \draw[arrow] (d) -- (psi);

    % Diramazioni ortogonali eleganti verso il box centrale eta(q, d)
    \draw[arrow] (q.south) -- ++(0, -0.5) -| ([xshift=-16pt]eta.north);
    \draw[arrow] (d.south) -- ++(0, -0.5) -| ([xshift=16pt]eta.north);

    % Connessioni dalle rappresentazioni alla funzione di aggregazione
    \draw[arrow] (phi.south) -- ++(0, -0.6) -| ([xshift=-40pt]f.north);
    \draw[arrow] (eta.south) -- (f.north);
    \draw[arrow] (psi.south) -- ++(0, -0.6) -| ([xshift=40pt]f.north);

    % Connessione dall'aggregatore allo score finale
    \draw[arrow] (f.south) -- (score.north);
\end{tikzpicture}
\caption{Decomposizione funzionale di una funzione di ranking generale basata sulle rappresentazioni (rielaborata da Guo et al.~\cite{guo2020}).}
\label{fig:ltr-general-ranking-framework}
\end{figure}

Nel modello Bag of Words classico:
\begin{itemize}
    \item La query e il documento sono proiettati come vettori sparsi di frequenze o pesi TF-IDF: $\phi(q) \in \mathbb{R}^{|\mathcal{V}|}$ e $\psi(d) \in \mathbb{R}^{|\mathcal{V}|}$;
    \item Il modulo di interazione esplicita \`e nullo: $\eta(q, d) = \emptyset$;
    \item La funzione di aggregazione $f$ consiste in una metrica chiusa prefissata, quale il coseno vettoriale o la sommatoria non lineare Okapi BM25:
    \begin{equation}
        s_{\text{BM25}}(q, d) = \sum_{t \in q \cap d} \text{IDF}(t) \cdot \frac{f_{t,d} \cdot (k_1 + 1)}{f_{t,d} + k_1 \cdot \left(1 - b + b \cdot \frac{|d|}{\text{avgdl}}\right)}.
    \end{equation}
\end{itemize}

Queste rappresentazioni sparse sono memorizzate negli indici invertiti, garantendo scansioni scalabili mediante algoritmi top-$k$ (WAND, MaxScore). Ciononostante, il modello BoW impone forti semplificazioni: assume l'indipendenza statistica tra i termini e trascura qualsiasi caratteristica strutturale del testo o dell'ambiente applicativo.

\subsection{Oltre il Bag of Words: La Molteplicit\`a dei Segnali di Rilevanza}
\label{subsec:segnali-oltre-bow}

Sul Web, una pagina non \`e un mero multinsieme non strutturato di vocaboli. Essa possiede una ricca anatomia interna e interagisce con un ecosistema complesso:
\begin{enumerate}
    \item \textbf{Campi e Zone Strutturate}: Un documento \`e suddiviso in campi eterogenei (titolo, abstract, corpo principale, testo dell'URL). L'occorrenza di un termine nel titolo trasmette un segnale di rilevanza drasticamente superiore rispetto all'occorrenza nel corpo.
    \item \textbf{Dati Testuali Esterni}: I testi di ancoraggio (\textit{anchor text}) degli hyperlink che puntano alla pagina rappresentano sintesi redatte da terzi, spesso pi\`u descrittive del contenuto stesso della risorsa.
    \item \textbf{Grafi di Collegamento e Topologia}: Link in entrata (\textit{in-links}), link in uscita (\textit{out-links}), score globali di autorit\`a quali PageRank e HITS.
    \item \textbf{Segnali Comportamentali e di Utilizzo}: Log delle ricerche, numero aggregato di clic ricevuti dall'URL per una determinata query, tempo di permanenza (\textit{dwell time}) sulla pagina prima di tornare ai risultati della SERP.
    \item \textbf{Propriet\`a della Query}: Lunghezza della query in termini, popolarit\`a nel log storico delle interrogazioni, entit\`a riconosciute, contesto geografico e temporale dell'utente.
\end{enumerate}

Di fronte a tale eterogeneit\`a, sorge spontanea la domanda fondante: \emph{Come possiamo apprendere automaticamente la funzione di aggregazione $f$ a partire dai dati anzich\'e concepirla manualmente?}

\FloatBarrier
\section{Il Paradigma del Machine Learning Applicato al Ranking}
\label{sec:paradigma-machine-learning-ranking}

Per superare le limitazioni delle formule euristiche chiuse, il problema del ranking viene riformulato all'interno del framework statistico del Machine Learning: consentire a un sistema computazionale di migliorare progressivamente le proprie prestazioni su un compito specifico attraverso l'esposizione all'esperienza (dati empirici) senza essere esplicitamente programmato.

\subsection{Formalizzazione dell'Apprendimento Supervisionato}
\label{subsec:formalizzazione-apprendimento-supervisionato}

Nel paradigma supervisionato, l'algoritmo riceve un insieme di istanze di addestramento (\textit{training instances}) descritte da un vettore di attributi $\mathbf{x} \in \mathcal{X}$ e corredate da un'etichetta desiderata $y \in \mathcal{Y}$ (\textit{ground truth} o supervisione fornita da annotatori umani):
\begin{itemize}
    \item Se $\mathcal{Y} \subseteq \mathbb{R}$, il task \`e di \textbf{Regressione} (stima di una grandezza continua);
    \item Se $\mathcal{Y} = \{c_1, \dots, c_m\}$ con $m \ge 2$, il task \`e di \textbf{Classificazione} (assegnazione a categorie discrete).
\end{itemize}

Il ciclo iterativo di addestramento si articola nei seguenti passi:
\begin{enumerate}
    \item Un'istanza $\mathbf{x}_i$ viene sottoposta al modello predittivo con parametri $\mathbf{\theta}$, generando una stima $\hat{y}_i = h_{\mathbf{\theta}}(\mathbf{x}_i)$;
    \item Viene quantificata la discrepanza tra la predizione $\hat{y}_i$ e l'etichetta vera $y_i$ attraverso una funzione di perdita $\mathcal{L}(\hat{y}_i, y_i)$ (\textit{loss function});
    \item I parametri $\mathbf{\theta}$ vengono aggiornati lungo la direzione che minimizza la perdita empirica (tipicamente mediante discesa del gradiente);
    \item Il procedimento viene ripetuto su tutti gli esempi del dataset per molteplici epoche di training (\textit{epochs}).
\end{enumerate}

\subsection{Analogia del Cervello e del Puzzle}
\label{subsec:analogia-puzzle}

Un'utile analogia concettuale per comprendere il ciclo retroattivo di ottimizzazione \`e quella della risoluzione di un rompicapo o puzzle:
\begin{itemize}
    \item A ogni passo abbiamo a disposizione il problema (\emph{il puzzle}) e la sua soluzione ideale certificata (\emph{ground truth});
    \item Il nostro modello cognitivo (\emph{il cervello / l'algoritmo}) genera un tentativo di soluzione;
    \item Confrontando la soluzione prodotta con quella corretta, rileviamo gli errori commessi;
    \item Sulla scorta dell'entit\`a e della direzione dello scarto, calibriamo i nostri parametri interni;
    \item Ripetendo la procedura per tentativi successivi su un vasto numero di puzzle, il sistema apprende una capacit\`a di generalizzazione efficace.
\end{itemize}

\subsection{Architettura a Due Fasi: Training Offline vs. Inference Online}
\label{subsec:architettura-offline-online}

Un vincolo imprescindibile nell'ingegneria dei motori di ricerca risiede nella rigorosa separazione temporale e computazionale tra le fasi di apprendimento e utilizzo del modello:

\begin{figure}[htbp]
\centering
\begin{tikzpicture}[
    phasebox/.style={rectangle, draw=navy, fill=codebackground, rounded corners=5pt, line width=0.9pt, inner sep=10pt, text width=0.88\linewidth, font=\small},
    arrow/.style={->, >=stealth, line width=1.2pt, crimson}
]
    \node[phasebox] (offline) {
        \textbf{\textcolor{navy}{\large FASE 1: TRAINING (OFFLINE)}}\\[4pt]
        $\bullet$ \textbf{Dataset Storico}: Insiemi di query $\mathcal{Q}_{\text{train}}$, collezioni di documenti candidati $\mathcal{D}$ ed etichette umane $y \in \{0, 1, \dots, 4\}$ (rilevanza graduata).\\[2pt]
        $\bullet$ \textbf{Estrazione Feature}: Calcolo esaustivo dei vettori $\mathbf{x} = \langle f_0, f_1, \dots, f_{K-1} \rangle$.\\[2pt]
        $\bullet$ \textbf{Algoritmo di ML}: Calcolo predizioni $\hat{y} = h_{\mathbf{\theta}}(\mathbf{x})$, valutazione della perdita $\mathcal{L}$, calcolo del gradiente $\nabla_{\mathbf{\theta}}\mathcal{L}$ e aggiornamento iterativo dei pesi $\mathbf{\theta}$.\\[2pt]
        $\bullet$ \textbf{Profilo Operativo}: Costo computazionale elevato (ore o giorni su cluster distribuiti), ma eseguito \emph{una tantum} fuori linea.
    };

    \node[phasebox, below=1.2cm of offline] (online) {
        \textbf{\textcolor{navy}{\large FASE 2: TESTING / INFERENCE (ONLINE)}}\\[4pt]
        $\bullet$ \textbf{Deployment}: Integrazione del modello validato $h_{\mathbf{\theta}^*}$ nella pipeline del search engine.\\[2pt]
        $\bullet$ \textbf{Query Real-Time}: Arrivo dell'interrogazione $q$ dall'utente in tempo reale.\\[2pt]
        $\bullet$ \textbf{Retrieval e Scoring}: Pre-selezione di $N$ candidati ed estrazione rapida delle feature.\\[2pt]
        $\bullet$ \textbf{Valutazione e Ordinamento}: Valutazione immediata di $h_{\mathbf{\theta}^*}(q, d)$ e ordinamento dei migliori $K$ risultati da restituire nella SERP.\\[2pt]
        $\bullet$ \textbf{Vincolo Ferreo}: Rispetto rigoroso del \textbf{Time Budget} di sistema ($\le 20\text{--}50\,\text{ms}$).
    };

    \draw[arrow] (offline) -- node[right, font=\small\bfseries, text=crimson, xshift=4pt] {Generazione del modello ottimizzato $\mathbf{\theta}^*$} (online);
\end{tikzpicture}
\caption{Disaccoppiamento strutturale tra la fase di addestramento offline e la fase di inferenza online.}
\label{fig:ltr-offline-online-phases}
\end{figure}

\FloatBarrier
\section{Definizione Formale di Learning to Rank e Tassonomia delle Feature}
\label{sec:definizione-ltr-tassonomia-feature}

Il framework denominato \textbf{Learning to Rank (LtR)} costituisce l'applicazione metodologica dell'apprendimento automatico al compito di costruire modelli di ordinamento per il recupero delle informazioni~\cite{liu2009}.

\subsection{Definizione e Obiettivo di Apprendimento}
\label{subsec:definizione-obiettivo-ltr}

\begin{definition}[Learning to Rank]
Sia data una collezione di query di training $\mathcal{Q}_{\text{train}} = \{q_1, q_2, \dots, q_m\}$. Per ogni query $q_i$, sia disponibile un insieme di documenti candidati $D_i = \{d_{i,1}, d_{i,2}, \dots, d_{i, n_i}\}$ e un corrispondente vettore di giudizi di rilevanza $\mathbf{y}_i = \langle y_{i,1}, y_{i,2}, \dots, y_{i, n_i} \rangle$, dove $y_{i,j} \in \mathcal{Y}$.  
Ciascuna coppia $(q_i, d_{i,j})$ \`e rappresentata da un vettore di attributi multidimensionale:
\begin{equation}
    \mathbf{x}_{i,j} = \Phi(q_i, d_{i,j}) = \langle f_0(q_i, d_{i,j}), f_1(q_i, d_{i,j}), \dots, f_{K-1}(q_i, d_{i,j}) \rangle \in \mathbb{R}^K.
    \label{eq:feature-vector-ltr}
\end{equation}
L'obiettivo di Learning to Rank consiste nell'apprendere una funzione di scoring $h: \mathbb{R}^K \to \mathbb{R}$ che assegni a ciascun documento un valore numerico tale per cui l'ordinamento indotto su $D_i$:
\begin{equation}
    \pi(D_i) = \text{sort}\Big(\{h(\mathbf{x}_{i,j})\}_{j=1}^{n_i}\Big)
\end{equation}
massimizzi una metrica globale di efficacia della graduatoria (come NDCG@$k$ o MAP).
\end{definition}

\begin{alertblock}{Distinzione Cruciale: Regressione vs Ranking}
Anche se i modelli LtR sono frequentemente addestrati come regressori o classificatori che assegnano un valore scalare al documento, l'obiettivo ultimo del sistema non \`e predire con precisione infinitesima la label numerica $y$, bens\`i \textbf{indovinare la permutazione corretta dell'intera lista}, massimizzando la rilevanza dei documenti collocati nelle posizioni di vertice.
\end{alertblock}

Le etichette di rilevanza umane $y$ sono tipicamente graduate su una scala discreta a $5$ livelli:
\begin{itemize}
    \item $y = 0$: \textit{Irrelevant} / Scadente;
    \item $y = 1$: \textit{Marginally Relevant} / Marginalmente utile;
    \item $y = 2$: \textit{Relevant} / Rilevante;
    \item $y = 3$: \textit{Highly Relevant} / Altamente rilevante;
    \item $y = 4$: \textit{Perfectly Relevant} / Risposta definitiva (ad es. home page cercata per query navigazionale).
\end{itemize}

\subsection{Tassonomia delle Feature in un Sistema di Ricerca Web}
\label{subsec:tassonomia-feature-ltr}

Nei motori di ricerca commerciali, la cardinalit\`a dello spazio delle feature \`e rilevante: gi\`a nel 2008 il \emph{New York Times} riportava le dichiarazioni di Amit Singhal secondo cui Google impiegava oltre $200$ feature distinte per ordinare le pagine web.

Secondo la classificazione canonica (Manning et al.~\cite{manning2008}, Liu~\cite{liu2009}), le feature si articolano in tre categorie strutturali distinte:

\begin{enumerate}
    \item \textbf{Feature Query-Only (Dipendenti esclusivamente dalla Query)}:
    Assumono il medesimo valore per tutti i documenti candidati associati alla stessa query. Non alterano l'ordinamento relativo dei documenti se applicate linearmente, ma fungono da selettori di contesto o modificatori di pendenza nei modelli ad albero:
    \begin{itemize}
        \item Numero di termini della query ($|q|$);
        \item Frequenza storica della query nei registri di log (popolarit\`a);
        \item Tipologia o intento della query (navigazionale, informativa, transazionale);
        \item Lingua rilevata o localizzazione geografica dell'interrogazione.
    \end{itemize}

    \item \textbf{Feature Query-Independent (Dipendenti esclusivamente dal Documento / statiche)}:
    Quantificano la qualit\`a intrinseca, l'autorevolezza o l'accessibilit\`a a priori del documento $d$, a prescindere da quale query lo abbia evocato:
    \begin{itemize}
        \item Valore di PageRank della pagina;
        \item Lunghezza dell'URL (in caratteri o in livelli di sottodirectory);
        \item Presenza di caratteri specifici nell'URL (es. tilde `\textasciitilde{}', indicativa di home page personali spesso obsolete);
        \item Numero di hyperlink in entrata (\textit{in-links}) e in uscita (\textit{out-links});
        \item Freschezza del documento (data dell'ultima modifica o tempo trascorso dall'indicizzazione);
        \item Lunghezza complessiva del testo del documento ($|d|$);
        \item Numero di immagini o elementi multimediali presenti nella pagina;
        \item Clic aggregati storici ricevuti dall'URL e Dwell Time medio globale.
    \end{itemize}

    \item \textbf{Feature Query-Dependent (Interazione Query--Documento)}:
    Misurano la rilevanza tematica e l'attinenza semantica specifica tra la query $q$ e il documento $d$:
    \begin{itemize}
        \item Punteggio Okapi BM25 calcolato sui diversi campi del documento;
        \item Similarit\`a del coseno nello spazio vettoriale VSM;
        \item Frequenza logaritmica dei termini della query presenti negli \textit{anchor text} dei link puntanti a $d$;
        \item Termini della query evidenziati in grassetto, con colori specifici o con tag tipografici primari (\texttt{<h1>}, \texttt{<title>});
        \item Finestra minima di prossimit\`a (\textit{term proximity}): la lunghezza della sequenza minima di testo che racchiude tutti i termini dell'interrogazione;
        \item Conteggio dei clic storici registrati per la specifica coppia $(q, d)$.
    \end{itemize}
\end{enumerate}

\subsection{L'Architettura di Pipeline Two-Stage e il Vincolo del Time Budget}
\label{subsec:pipeline-two-stage-budget}

Come evidenziato nella Figura~\ref{fig:ltr-two-stage-pipeline}, un motore di ricerca moderno adotta un'architettura a imbuto (\textit{Two-Stage Retrieval}):
\begin{itemize}
    \item \textbf{Primo Stadio (Base Ranker)}: Interroga l'indice invertito compresso mediante algoritmi ad alta efficienza (come WAND o MaxScore descritti nel Capitolo~\ref{chap:elaborazione-topk-wand-maxscore}). Valuta un punteggio lessicale economico (BM25) ed estrae una rosa ristretta di $N \approx 1\,000$ documenti candidati.
    \item \textbf{Secondo Stadio (Top Ranker / LtR Ranker)}: Riceve gli $N$ candidati, calcola i vettori di feature multidimensionali $\mathbf{x} \in \mathbb{R}^K$ ed esegue il modello di Machine Learning appreso offline per riordinare i documenti e selezionare i migliori $k \approx 10\text{--}20$ risultati da presentare all'utente nella SERP.
\end{itemize}

\begin{figure}[htbp]
\centering
\resizebox{\linewidth}{!}{%
\begin{tikzpicture}[
    basebox/.style={
        rectangle, 
        draw=navy, 
        fill=codebackground, 
        rounded corners=4pt, 
        line width=0.8pt, 
        align=center, 
        inner sep=6pt, 
        font=\small
    },
    procbox/.style={
        rectangle, 
        draw=darkgreen, 
        fill=darkgreen!10, 
        rounded corners=4pt, 
        line width=0.9pt, 
        align=center, 
        inner sep=6pt, 
        font=\small\bfseries
    },
    resbox/.style={
        rectangle, 
        draw=crimson, 
        fill=crimson!10, 
        rounded corners=4pt, 
        line width=0.9pt, 
        align=center, 
        inner sep=6pt, 
        font=\small\bfseries
    },
    auxbox/.style={
        rectangle, 
        draw=navy!70, 
        fill=navy!5, 
        dashed,
        rounded corners=4pt, 
        line width=0.8pt, 
        align=center, 
        inner sep=5pt, 
        font=\footnotesize
    },
    arrow/.style={->, >=stealth, line width=1.0pt, navy},
    dasharrow/.style={->, >=stealth, line width=0.9pt, dashed, navy!80}
]
    % --- FLUSSO ORIZZONTALE PRINCIPALE (y = 0) ---
    % Query (x = 0)
    \node[basebox, text width=2.2cm, minimum height=1.2cm] (query) at (0, 0) {Query Utente\\[2pt]$q$};

    % Primo Stadio (x = 3.6)
    \node[procbox, text width=3.0cm, minimum height=1.2cm] (first) at (3.6, 0) {Primo Stadio\\[2pt]\footnotesize(Base Ranker, WAND)};

    % Candidati intermedi (x = 7.5)
    \node[basebox, text width=2.8cm, minimum height=1.2cm] (ncands) at (7.5, 0) {$N$ Documenti\\[2pt]Candidati\\[2pt]\footnotesize($N \approx 1\,000$)};

    % Secondo Stadio (x = 11.8)
    \node[procbox, text width=3.4cm, minimum height=1.2cm] (second) at (11.8, 0) {Secondo Stadio\\[2pt]\footnotesize(Top Ranker / LtR)};

    % Risultati SERP (x = 15.8)
    \node[resbox, text width=2.8cm, minimum height=1.2cm] (serp) at (15.8, 0) {Risultati SERP\\[2pt]Top-$k$\\[2pt]\footnotesize($k \approx 10\text{--}20$)};

    % --- NODI AUSILIARI SUPERIORI (y = 2.0) ---
    % Indice Invertito sopra al Primo Stadio
    \node[auxbox, text width=2.8cm] (index) at (3.6, 2.0) {Indice Invertito\\dei Documenti};

    % Estrazione Feature sopra al Secondo Stadio
    \node[auxbox, text width=3.0cm] (feats) at (11.8, 2.0) {Estrazione Feature\\$\mathbf{x} = \Phi(q, d) \in \mathbb{R}^K$};

    % --- CONNESSIONI DEL FLUSSO PRINCIPALE ---
    \draw[arrow] (query) -- (first);
    \draw[arrow] (first) -- (ncands);
    \draw[arrow] (ncands) -- (second);
    \draw[arrow] (second) -- (serp);

    % --- CONNESSIONI AUSILIARIE ---
    % Da Indice Invertito a Primo Stadio (discesa verticale)
    \draw[dasharrow] (index.south) -- (first.north);

    % Da Candidati verso l'alto verso Estrazione Feature e poi nel Secondo Stadio
    \draw[dasharrow] (ncands.north) |- (feats.west);
    \draw[dasharrow] (feats.south) -- (second.north);

\end{tikzpicture}%
}
\caption{Pipeline di elaborazione a due stadi: dal filtraggio scalabile su indice invertito al re-ranking supervisionato.}
\label{fig:ltr-two-stage-pipeline}
\end{figure}

Il costo computazionale dei modelli LtR influenza direttamente il throughput complessivo dell'infrastruttura di servizio. I vincoli di latenza impongono che l'estrazione delle centinaia di feature e la valutazione del modello predittivo avvengano entro una frazione infinitesima del budget temporale complessivo ($\le 15\text{--}20\,\text{ms}$ per query). Per questa ragione, la ricerca recente ha sviluppato ottimizzazioni architetturali di basso livello (implementazioni tree-traversal branchless, compilazione di foreste decisionali in codice assembly nativo) volte ad accelerare l'inferenza online.

\FloatBarrier
\section{Classificazione per l'Ad-Hoc IR: Formulazione ed Esempio Geometrico}
\label{sec:classificazione-adhoc-ir-geometria}

Uno dei primi approcci all'integrazione del Machine Learning nel ranking ha tentato di ridurre la selezione dei documenti a un problema classico di \textbf{classificazione lineare binaria} (Manning et al.~\cite{manning2008}).

\subsection{Formulazione Lineare dello Scoring}
\label{subsec:formulazione-lineare-scoring}

Si consideri un dataset di addestramento composto da triple $(q, d, r)$, dove la rilevanza $r \in \{0, 1\}$ \`e binaria: $r = 1$ identifica la classe dei documenti rilevanti ($R$), mentre $r = 0$ identifica la classe dei non rilevanti ($N$).

Ciascuna coppia $(q, d)$ viene proiettata in un semplice spazio bidimensionale $\mathbf{x} = \langle \alpha, \omega \rangle$:
\begin{itemize}
    \item $\alpha \in [0, 1]$: la similarit\`a del coseno tra la query e il documento calcolata secondo il modello VSM;
    \item $\omega \ge 1$: la minima finestra di prossimit\`a (\textit{minimum query window size}), definita come la lunghezza del pi\`u breve segmento contiguo di testo che include tutti i termini della query (indipendentemente dall'ordine). Una finestra $\omega$ ridotta indica che i termini compaiono ravvicinati, fornendo una prova tangibile di coerenza semantica.
\end{itemize}

Si definisce una funzione di punteggio lineare parametrizzata dai coefficienti scalari $a$ e $b$:
\begin{equation}
    \text{Score}(q, d) = \text{Score}(\langle \alpha, \omega \rangle) = a \cdot \alpha + b \cdot \omega.
    \label{eq:linear-score-function}
\end{equation}
Il classificatore lineare stabilisce una superficie di decisione attraverso una soglia critica $\Theta \in \mathbb{R}$:
\begin{equation}
    \text{Classe}(q, d) = 
    \begin{cases}
        \text{Rilevante } (R), & \text{se } \text{Score}(q, d) > \Theta, \\
        \text{Non Rilevante } (N), & \text{se } \text{Score}(q, d) \le \Theta.
    \end{cases}
    \label{eq:linear-decision-rule}
\end{equation}
L'apprendimento consiste nell'individuare la tupla ottimale di parametri $\langle a^*, b^*, \Theta^* \rangle$ in grado di separare le due classi nello spazio geometrico minimizzando l'errore empirico di classificazione.

\subsection{Esempio Geometrico e Superficie di Decisione}
\label{subsec:esempio-geometrico-superficie}

Nella Figura~\ref{fig:ltr-linear-classification-space} viene visualizzato lo spazio bidimensionale $(\omega, \alpha)$ con la disposizione empirica delle osservazioni e la retta di separazione:

\begin{figure}[htbp]
\centering
\begin{tikzpicture}[
    >=Stealth,
    decisionline/.style={dashed, line width=1.1pt, navy},
    markerR/.style={circle, fill=darkgreen!85!black, draw=darkgreen, line width=0.6pt, inner sep=2.2pt, opacity=0.9},
    markerN/.style={regular polygon, regular polygon sides=3, fill=crimson!85!black, draw=crimson, line width=0.6pt, inner sep=2.0pt, opacity=0.9}
]

    % Margini di respiro laterali per evitare troncamenti a sinistra e a destra
    \path (-1.4, 0) (8.9, 0);

    % ==================================================================
    % 0. CAMPITURE DEI SEMIPIANI DECISIONALI (Verde e Rosso pastello)
    % Retta: y = 0.5*x + 1.35 da x=0 a x=7.5 (y=5.10)
    % ==================================================================
    \begin{scope}[on background layer]
        % Regione Predetta Rilevante (R) - Verde chiarissimo
        \fill[darkgreen!8] (0, 1.35) -- (7.5, 5.10) -- (7.5, 5.3) -- (0, 5.3) -- cycle;
        % Regione Predetta Non Rilevante (N) - Rosso chiarissimo
        \fill[crimson!5] (0, 0) -- (7.5, 0) -- (7.5, 5.10) -- (0, 1.35) -- cycle;

        % Griglia cartesiana ortogonale regolare (passo 0.01 su y, unitario su x)
        \foreach \x in {0, 1.5, 3.0, 4.5, 6.0} {
            \draw[gray!20, dashed, line width=0.4pt] (\x, 0) -- (\x, 5.3);
        }
        \foreach \y in {1.0, 2.0, 3.0, 4.0, 5.0} {
            \draw[gray!20, dashed, line width=0.4pt] (0, \y) -- (7.5, \y);
        }
    \end{scope}

    % ==================================================================
    % 1. ASSI CARTESIANI E SCALE REGOLARI
    % ==================================================================
    \draw[->, line width=1.0pt, navy] (0, 0) -- (7.8, 0)
        node[right=3pt, font=\footnotesize\bfseries, text=navy] {$\omega$};
    \draw[->, line width=1.0pt, navy] (0, 0) -- (0, 5.4)
        node[above=2pt, font=\footnotesize\bfseries, text=navy] {$\alpha$};

    % Etichette assi distanziate in modo ottimale per la massima leggibilità
    \node[below=22pt, font=\small\bfseries, text=navy] at (3.75, 0) {Term proximity $\omega$ (ampiezza minima finestra)};
    \node[rotate=90, anchor=south, font=\small\bfseries, text=navy] at (-1.05, 2.5) {Cosine score $\alpha$};

    % Tacche asse Y regolari a passo 0.01 (0, 0.01, 0.02, 0.03, 0.04, 0.05)
    \foreach \y/\val in {0/0, 1.0/0.01, 2.0/0.02, 3.0/0.03, 4.0/0.04, 5.0/0.05} {
        \draw[navy, line width=0.6pt] (0, \y) -- ++(-3pt, 0) node[left=3pt, font=\scriptsize] {$\val$};
    }

    % Tacche asse X su omega (2, 3, 4, 5, 6)
    \foreach \x/\val in {0/2, 1.5/3, 3.0/4, 4.5/5, 6.0/6} {
        \draw[navy, line width=0.6pt] (\x, 0) -- ++(0, -3pt) node[below=3pt, font=\scriptsize] {$\val$};
    }

    % ==================================================================
    % 2. SUPERFICIE DI DECISIONE LINEARE (Retta pulita con etichetta terminale)
    % ==================================================================
    \draw[decisionline] (0, 1.35) -- (7.5, 5.10)
        node[anchor=south west, font=\footnotesize\bfseries, text=navy] {$a\alpha + b\omega = \Theta$};

    % ==================================================================
    % 3. CAMPIONI OSSERVATI (Marker geometrici differenziati con lieve jittering)
    % ==================================================================
    % Colonna \omega = 2 (x = 0.12)
    \node[markerR] at (0.12, 1.85) {};

    % Colonna \omega = 3 (x = 1.5)
    \node[markerR] at (1.50, 4.05) {};
    \node[markerR] at (1.46, 2.65) {};
    \node[markerR] at (1.54, 2.25) {};
    \node[markerN] at (1.48, 1.70) {};
    \node[markerN] at (1.52, 0.75) {};

    % Colonna \omega = 4 (x = 3.0)
    \node[markerR] at (3.00, 3.85) {};
    \node[markerR] at (3.04, 3.45) {};
    % Falso Negativo: R sotto la retta tratteggiata!
    \node[markerR] (fn) at (2.96, 2.35) {};
    \draw[crimson, line width=0.9pt, dashed] (fn) circle (5.5pt);
    \draw[->, line width=0.6pt, crimson] (fn) ++(6pt, 0) -- ++(12pt, 0) 
        node[right, font=\tiny\bfseries, text=crimson] {Falso Negativo ($R \to N$)};

    \node[markerN] at (3.03, 1.85) {};
    \node[markerN] at (2.97, 1.20) {};
    \node[markerN] at (3.02, 0.65) {};

    % Colonna \omega = 5 (x = 4.5)
    \node[markerR] at (4.50, 4.35) {};
    \node[markerR] at (4.46, 4.05) {};
    \node[markerR] at (4.54, 3.65) {};
    \node[markerN] at (4.48, 2.75) {};
    \node[markerN] at (4.52, 0.85) {};

    % Colonna \omega = 6 (x = 6.0)
    \node[markerR] at (6.00, 4.85) {};
    % Falso Positivo: N sopra la retta tratteggiata!
    \node[markerN] (fp) at (5.96, 4.40) {};
    \draw[crimson, line width=0.9pt, dashed] (fp) circle (5.5pt);
    % Freccia rivolta a destra verso lo spazio libero
    \draw[->, line width=0.6pt, crimson] (fp) ++(6pt, 0) -- ++(12pt, 0) 
        node[right, font=\tiny\bfseries, text=crimson] {Falso Positivo ($N \to R$)};

    \node[markerN] at (6.03, 2.75) {};
    \node[markerN] at (5.98, 1.65) {};

    % ==================================================================
    % 4. LEGENDA ACCADEMICA STRUTTURATA (Esterna agli assi cartesiani)
    % ==================================================================
    \node[draw=gray!35, fill=white, rounded corners=3pt, line width=0.5pt, anchor=south, inner sep=4.5pt] at (3.75, 6.3) {
        \scriptsize
        \begin{tabular}{c}
        \tikz\node[markerR, inner sep=2pt] {};~Rilevante ($R$) \qquad
        \tikz\node[markerN, inner sep=1.8pt] {};~Non Rilevante ($N$) \qquad
        \tikz\draw[dashed, line width=0.8pt, crimson] (0, 0) circle (3pt);~Errore empirico \\[3pt]
        \tikz\fill[darkgreen!18, draw=darkgreen!60, line width=0.4pt, rounded corners=1pt] (0,-0.08) rectangle (0.3,0.15);~Regione predetta $R$ ($a\alpha + b\omega > \Theta$) \qquad\qquad
        \tikz\fill[crimson!12, draw=crimson!60, line width=0.4pt, rounded corners=1pt] (0,-0.08) rectangle (0.3,0.15);~Regione predetta $N$ ($a\alpha + b\omega \le \Theta$)
        \end{tabular}
    };

\end{tikzpicture}
\caption{Spazio delle feature $\langle\omega, \alpha\rangle$ e classificazione lineare per Information Retrieval (Manning et al.~\cite{manning2008}, Sez.~15.4.1): cerchi verdi per documenti rilevanti ($R$), triangoli rossi per non rilevanti ($N$). La retta tratteggiata $a\alpha + b\omega = \Theta$ modella il confine di decisione lineare. Le aree ombreggiate identificano le due semiregioni decisionali previste; i cerchi tratteggiati evidenziano il Falso Negativo su $\omega=4$ e il Falso Positivo su $\omega=6$, a dimostrazione che il dataset empirico non \`e linearmente separabile al 100\%.}
\label{fig:ltr-linear-classification-space}
\end{figure}

Il grafico visualizza la distribuzione empirica dei documenti nello spazio bidimensionale definito dalla prossimit\`a dei termini ($\omega$) sull'asse orizzontale e dal punteggio di similarit\`a coseno ($\alpha$) sull'asse verticale. La retta tratteggiata $a\alpha + b\omega = \Theta$ modella il confine di decisione lineare: all'aumentare della distanza tra i termini (peggioramento qualitativo di $\omega$), un documento necessita di un punteggio di coseno $\alpha$ progressivamente pi\`u elevato per essere classificato come rilevante ($R$). Le aree ombreggiate identificano le due semiregioni decisionali previste dal modello; la presenza di campioni oltre la soglia lineare evidenzia come il dataset empirico non sia perfettamente separabile in modo lineare.

Dall'analisi geometrica e concettuale emergono propriet\`a fondamentali per l'Information Retrieval:

\begin{enumerate}
    \item \textbf{Significato Fisico degli Assi e Regione di Rilevanza}:
    \begin{itemize}
        \item L'asse delle ordinate misura il punteggio coseno $\alpha \in [0, 1]$: valori pi\`u alti corrispondono a una maggiore similarit\`a globale tra la query e il documento.
        \item L'asse delle ascisse misura la dimensione minima della finestra di testo $\omega \ge 2$ che racchiude tutti i termini della query nel documento: trattandosi di una distanza, valori inferiori riflettono una maggiore prossimit\`a sintattica e concettuale (quindi migliore qualit\`a).
        \item La regione di massima rilevanza si concentra pertanto in \textbf{alto a sinistra} (elevato coseno $\alpha$ e ridotta finestra $\omega$).
    \end{itemize}

    \item \textbf{Pendenza Positiva della Retta di Decisione (Trade-Off Compensativo)}:
    \begin{itemize}
        \item La superficie di separazione lineare $a\alpha + b\omega = \Theta$ esibisce pendenza positiva nel piano cartesiano.
        \item Questa geometria modella formalmente un principio di \emph{compensazione}: all'aumentare di $\omega$ (ovvero quando i termini della query si disperdono all'interno del documento, degradandone la prossimit\`a), il documento pu\`o comunque essere classificato come rilevante ($R$) solo se compensa tale carenza con un punteggio di coseno $\alpha$ proporzionalmente pi\`u elevato.
    \end{itemize}

    \item \textbf{Non-Separabilit\`a Lineare Perfetta ed Errori Empirici}:
    \begin{itemize}
        \item I campioni reali non risultano linearmente separabili al 100\%: l'iperpiano di decisione ottimale minimizza l'errore empirico, ma non pu\`o annullarlo.
        \item \textbf{Falso Negativo} ($R$ sotto la soglia): sulla colonna $\omega = 4$, il campione rilevante a quota $\alpha \approx 0.024$ si colloca al di sotto della retta, venendo classificato erroneamente come non rilevante ($N$).
        \item \textbf{Falso Positivo} ($N$ sopra la soglia): sulla colonna $\omega = 6$, il campione non rilevante a quota $\alpha \approx 0.044$ si trova al di sopra della linea di decisione, venendo predetto erroneamente come rilevante ($R$).
    \end{itemize}
\end{enumerate}

Sebbene il classificatore lineare fornisca un'intuitiva demarcazione dello spazio, trattare il ranking documentale come un problema di pura classificazione binaria presenta limiti intrinseci severi:
\begin{itemize}
    \item \textbf{Assenza di ranking interno}: tutti i documenti che cadono al di sopra della retta appartengono indistintamente alla classe $R$, senza fornire un ordinamento di preferenza relativo (ad esempio tra il primo e il decimo risultato);
    \item \textbf{Simmetria ingannevole dei costi di errore}: una misclassificazione a quota altissima (documento cardine da prima posizione scartato) viene penalizzata dalla loss di classificazione esattamente con lo stesso peso di una misclassificazione su documenti a bassa rilevanza;
    \item \textbf{Incapacit\`a di gestire la rilevanza graduata}: le collezioni moderne impiegano scale multi-livello (es.~$\{0, 1, 2, 3, 4\}$ per NDCG), impossibili da catturare fedelmente con un singolo iperpiano binario.
\end{itemize}

\FloatBarrier
\section{La Tassonomia Classica di Learning to Rank: Pointwise, Pairwise e Listwise}
\label{sec:tassonomia-ltr-pointwise-pairwise-listwise}

Per colmare le discrepanze tra la formulazione classica del Machine Learning e le esigenze del reperimento di documenti, la letteratura scientifica ha codificato tre paradigmi fondamentali~\cite{liu2009}:

\begin{figure}[htbp]
\centering
\resizebox{\linewidth}{!}{%
\begin{tikzpicture}[
    font=\small\sffamily
]

    % Larghezza standard delle colonne (4.8 cm per evitare sillabazioni forzate)
    \def\cardW{4.8cm}

    % ==================================================================
    % 1. CORPI TESTUALI (Posizionati a y = 0, sviluppo verso il basso)
    % ==================================================================
    \node[text width=\cardW, align=flush left, inner sep=7pt, anchor=north] (body1) at (0, 0) {
        \textbf{Input:} Istanza singola\\
        $(\mathbf{x}(q, d),\, y)$.\\[5pt]
        \textbf{Target:} Stima dello score o classe di rilevanza $y \in \mathcal{Y}$.\\[5pt]
        \textbf{Modello:} Regressione standard o classificazione multiclasse.\\[5pt]
        \textbf{Criticit\`a:} Indifferente all'ordine relativo e alla posizione finale.
    };

    \node[text width=\cardW, align=flush left, inner sep=7pt, anchor=north] (body2) at (5.5, 0) {
        \textbf{Input:} Coppia di documenti per query: $(\mathbf{x}_i, \mathbf{x}_j,\, y_{ij})$.\\[5pt]
        \textbf{Target:} Preservazione dell'ordine relativo ($d_i \succ d_j$).\\[5pt]
        \textbf{Modello:} Classificazione binaria su coppie (es.~RankNet).\\[5pt]
        \textbf{Criticit\`a:} Pesa ogni coppia in modo identico, ignorando il vertice top-$k$.
    };

    \node[text width=\cardW, align=flush left, inner sep=7pt, anchor=north] (body3) at (11.0, 0) {
        \textbf{Input:} Insieme ordinato dei documenti associati a $q$.\\[5pt]
        \textbf{Target:} Ottimizzazione di metriche globali (NDCG, MAP).\\[5pt]
        \textbf{Modello:} Metodi a loss surrogata differenziabile (es.~LambdaMART).\\[5pt]
        \textbf{Criticit\`a:} Metriche IR discrete e non derivabili sui parametri.
    };

    % ==================================================================
    % 2. CALCOLO ALTEZZA COMUNE (Fondo allineato alla colonna più lunga)
    % ==================================================================
    \path let \p1 = (body1.south), \p2 = (body2.south), \p3 = (body3.south),
              \n1 = {min(\y1, min(\y2, \y3)) - 6pt} in
          coordinate (southCommon) at (0, \n1);

    % ==================================================================
    % 3. DISEGNO DELLE CARD (Cornice arrotondata + Banner superiore)
    % ==================================================================
    \begin{scope}[on background layer]
        \foreach \i/\title in {1/Pointwise, 2/Pairwise, 3/Listwise} {
            % Coordinate dei quattro vertici della card
            \path (body\i.north west) ++(-2pt, 0) coordinate (tl\i)
                  (body\i.north east) ++(2pt, 0) coordinate (tr\i)
                  (tl\i |- southCommon) coordinate (bl\i)
                  (tr\i |- southCommon) coordinate (br\i);

            % Sfondo con angoli arrotondati
            \draw[draw=navy, fill=codebackground, line width=0.9pt, rounded corners=5pt]
                ([yshift=26pt]tl\i) rectangle (br\i);

            % Maschera per arrotondare solo gli angoli superiori della barra navy
            \begin{scope}
                \clip[rounded corners=5pt]
                    ([yshift=26pt]tl\i) rectangle (br\i);
                \fill[navy]
                    ([yshift=27pt]tl\i) rectangle ([yshift=0pt]tr\i);
            \end{scope}

            % Ridisegno del bordo
            \draw[draw=navy, line width=0.9pt, rounded corners=5pt]
                ([yshift=26pt]tl\i) rectangle (br\i);
        }
    \end{scope}

    % Titoli centrati nelle testate
    \foreach \i/\title in {1/Pointwise, 2/Pairwise, 3/Listwise} {
        \node[font=\normalsize\bfseries, text=white, anchor=center] 
            at ($([yshift=13pt]tl\i)!0.5!([yshift=13pt]tr\i)$) {\title};
    }

\end{tikzpicture}%
}
\caption{Confronto sistematico dei tre approcci canonici di Learning to Rank: \mbox{granularit\`a} di input, target e relative criticit\`a.}
\label{fig:ltr-three-approaches-comparison}
\end{figure}

\subsection{Approccio Pointwise}
\label{subsec:approccio-pointwise}

Nel paradigma \textbf{Pointwise}:
\begin{itemize}
    \item Ogni coppia query-documento \`e associata a uno score indipendente $y \in \mathcal{Y}$;
    \item Il problema \`e formulato come regressione convenzionale (minimizzazione dell'errore quadratico medio $\mathcal{L}_2 = (h(\mathbf{x}) - y)^2$) oppure come classificazione multiclasse con loss di entropia incrociata;
    \item \textbf{Criticit\`a Fondamentale}: L'approccio \`e cieco rispetto alla struttura della graduatoria. Ad esempio, penalizza allo stesso modo una discrepanza di punteggio su un documento irrilevante a fondo lista e su un documento cardine collocato in prima posizione.
\end{itemize}

\subsection{Approccio Pairwise e il Modello RankNet}
\label{subsec:approccio-pairwise-ranknet}

Il paradigma \textbf{Pairwise} sposta l'unit\`a di analisi dalla stima assoluta alla preservazione dell'ordinamento relativo tra coppie di documenti associati alla medesima query:
\begin{itemize}
    \item Dati due documenti $d_i$ e $d_j$ candidati per la query $q$, se $y_i > y_j$ si stabilisce una preferenza ordinata $d_i \succ d_j$ (documento $d_i$ pi\`u rilevante di $d_j$);
    \item Il compito si riconduce a una classificazione binaria: il modello deve assegnare punteggi tali per cui $h(\mathbf{x}_i) > h(\mathbf{x}_j)$;
    \item La perdita empirica quantifica il numero di inversioni o violazioni di coppia (\emph{pairwise ranking loss}).
\end{itemize}

\subsubsection{Formulazione Algebrica di RankNet}
L'algoritmo \textbf{RankNet} (Burges et al.~\cite{burges2010}) modella probabilisticamente la preferenza di coppia tramite una funzione logistica sigmoide:
\begin{equation}
    P_{ij} = P(d_i \succ d_j) = \frac{1}{1 + e^{-\sigma (s_i - s_j)}},
    \label{eq:ranknet-probability}
\end{equation}
dove $s_i = h(\mathbf{x}_i)$ e $s_j = h(\mathbf{x}_j)$ denotano i punteggi predetti dal modello (ad esempio una rete neurale o un ensemble), e $\sigma > 0$ governa la pendenza della curva sigmoide.

Definita la probabilit\`a target della ground truth come $\bar{P}_{ij} = 1$ se $y_i > y_j$ (oppure $0.5$ se $y_i = y_j$), la funzione di costo di RankNet per la coppia $(d_i, d_j)$ corrisponde alla cross-entropy:
\begin{equation}
    \mathcal{C}_{ij} = -\bar{P}_{ij} \log P_{ij} - (1 - \bar{P}_{ij}) \log(1 - P_{ij}).
\end{equation}
Nel caso rigoroso in cui $d_i \succ d_j$ ($\bar{P}_{ij} = 1$), la funzione di costo si semplifica in forma analitica chiusa:
\begin{equation}
    \mathcal{C}_{ij} = \log\big(1 + e^{-\sigma (s_i - s_j)}\big).
    \label{eq:ranknet-loss-pair}
\end{equation}
La formulazione~\eqref{eq:ranknet-loss-pair} \`e ovunque continua e derivabile rispetto ai punteggi $s_i$ e $s_j$:
\begin{equation}
    \frac{\partial \mathcal{C}_{ij}}{\partial s_i} = -\sigma \left(\frac{1}{1 + e^{\sigma (s_i - s_j)}}\right) = -\sigma (1 - P_{ij}) = -\frac{\partial \mathcal{C}_{ij}}{\partial s_j}.
    \label{eq:ranknet-derivatives}
\end{equation}
Grazie alla differenziabilit\`a, RankNet consente l'addestramento diretto dei pesi interni tramite Stochastic Gradient Descent (SGD) applicando la regola della catena:
\begin{equation}
    \frac{\partial \mathcal{C}_{ij}}{\partial \mathbf{\theta}} = \frac{\partial \mathcal{C}_{ij}}{\partial s_i} \frac{\partial s_i}{\partial \mathbf{\theta}} + \frac{\partial \mathcal{C}_{ij}}{\partial s_j} \frac{\partial s_j}{\partial \mathbf{\theta}}.
\end{equation}

\subsection{Criticit\`a Strutturale del Pairwise e Divergenza con NDCG}
\label{subsec:criticita-pairwise-ndcg}

Nonostante l'eleganza matematica di RankNet, il paradigma Pairwise soffre di una patologia intrinseca: minimizzare il numero globale di coppie invertite non equivale a massimizzare metriche IR posizionali quali NDCG@$k$ o MAP.

\begin{alertblock}{Il Dilemma della Priorit\`a di Posizione}
Nella metrica Normalized Discounted Cumulative Gain (NDCG@$k$), il contributo di ciascun documento \`e ponderato da uno smorzamento logaritmico della posizione: $\frac{1}{\log_2(\text{rank} + 1)}$.  
Ne consegue che un'inversione tra la posizione $1$ e la posizione $2$ provoca una perdita di qualit\`a drasticamente superiore rispetto alla medesima inversione tra la posizione $80$ e la posizione $81$. L'approccio Pairwise classico tratta invece entrambe le inversioni come una singola violazione identica, distorcendo l'ottimizzazione del modello.
\end{alertblock}

Questa divergenza \`e documentata graficamente nella Figura~\ref{fig:ltr-pairwise-loss-vs-ndcg}, dove si osserva l'andamento della Pairwise Loss a confronto con NDCG@5 al progredire delle epoche di addestramento:

\begin{figure}[htbp]
\centering
\begin{tikzpicture}[
    ndcgline/.style={line width=1.3pt, crimson},
    lossline/.style={line width=1.3pt, teal!80!navy}
]
    % --- RETICOLO DI SFONDO (Tratteggiato sottile e allineato alle tacche comuni) ---
    \begin{scope}[on background layer]
        % Griglia verticale sulle epoche
        \foreach \x in {1.0, 2.0, 3.0, 4.0, 5.0, 6.0, 7.0, 8.0} {
            \draw[gray!20, dashed, line width=0.4pt] (\x, 0) -- (\x, 4.5);
        }
        % Griglia orizzontale: coincide esattamente sia sui valori di sinistra che di destra!
        % y = 1.5 -> NDCG = 0.48, Loss = 0.30
        % y = 3.0 -> NDCG = 0.49, Loss = 0.32
        \foreach \y in {1.5, 3.0} {
            \draw[gray!20, dashed, line width=0.4pt] (0, \y) -- (9.0, \y);
        }
    \end{scope}

    % --- ASSI ORIZZONTALI (Superiore e Inferiore in Navy neutro) ---
    \draw[navy, line width=0.8pt] (0, 0) -- (9.0, 0);
    \draw[navy, line width=0.8pt] (0, 4.5) -- (9.0, 4.5);

    % Tacche asse orizzontale: Epoche da 20 a 200
    \foreach \x/\val in {0/20, 1/40, 2/60, 3/80, 4/100, 5/120, 6/140, 7/160, 8/180, 9/200} {
        \draw[navy, line width=0.6pt] (\x, 0) -- ++(0, -3pt) node[below=2pt, font=\scriptsize] {\val};
        \draw[navy, line width=0.6pt] (\x, 4.5) -- ++(0, 3pt);
    }
    \node[font=\footnotesize\bfseries, text=navy] at (4.5, -0.65) {Epoche di addestramento};

    % --- ASSE Y SINISTRO: METRICA NDCG@5 (Codice Colore Crimson/Bordeaux) ---
    \draw[crimson, line width=1.0pt] (0, 0) -- (0, 4.5);
    \foreach \y/\val in {0.0/0.47, 1.5/0.48, 3.0/0.49, 4.5/0.50} {
        \draw[crimson, line width=0.7pt] (0, \y) -- ++(-3pt, 0) node[left=3pt, font=\scriptsize\bfseries, text=crimson] {$\val$};
    }
    \node[rotate=90, font=\footnotesize\bfseries, text=crimson] at (-1.1, 2.25) {Metrica posizionale (NDCG@5)};

    % --- ASSE Y DESTRO: FUNZIONE DI PERDITA (Codice Colore Ottanio/Teal) ---
    \draw[teal!80!navy, line width=1.0pt] (9.0, 0) -- (9.0, 4.5);
    \foreach \y/\val in {0.0/0.28, 1.5/0.30, 3.0/0.32, 4.5/0.34} {
        \draw[teal!80!navy, line width=0.7pt] (9.0, \y) -- ++(3pt, 0) node[right=3pt, font=\scriptsize\bfseries, text=teal!80!navy] {$\val$};
    }
    \node[rotate=-90, font=\footnotesize\bfseries, text=teal!80!navy] at (10.1, 2.25) {Funzione di perdita (Pairwise Loss)};

    % --- CURVA 1: PAIRWISE LOSS (Monotona decrescente da 0.334 a 0.285) ---
    % Mappatura: y = ((loss - 0.28) / 0.06) * 4.5
    \draw[lossline] plot[mark=x, mark size=1.6pt] coordinates {
        (0.0, 4.05) (0.5, 3.75) (1.0, 3.42) (1.5, 3.10) (2.0, 2.76)
        (2.5, 2.40) (3.0, 2.05) (3.5, 1.70) (4.0, 1.38) (4.5, 1.15)
        (5.0, 0.95) (5.5, 0.72) (6.0, 0.52) (6.5, 0.35) (7.0, 0.33)
        (7.5, 0.32) (8.0, 0.31) (8.5, 0.30) (9.0, 0.29)
    };

    % --- CURVA 2: NDCG@5 (Taylor et al. CIKM 2006) ---
    % Sale rapida: da 0.477 a 0.493 (epoca 50), poi degrado progressivo e stabilizzazione
    % Mappatura: y = ((ndcg - 0.47) / 0.03) * 4.5
    \draw[ndcgline] plot[mark=+, mark size=1.8pt] coordinates {
        (0.0, 1.05) (0.4, 1.55) (0.8, 2.20) (1.2, 3.45) (1.5, 3.30)
        (1.8, 3.15) (2.1, 3.15) (2.4, 3.55) (2.7, 2.40) (3.0, 2.30)
        (3.4, 2.45) (3.8, 1.95) (4.2, 1.50) (4.6, 1.45) (5.0, 1.10)
        (5.3, 1.35) (5.6, 0.85) (6.0, 0.80) (6.3, 0.75) (6.5, 1.25)
        (7.0, 1.25) (7.5, 1.25) (8.0, 1.25) (8.5, 1.25) (9.0, 1.25)
    };

    % --- LEGENDA ACCADEMICA COMPATTA (In alto, centrata e ariosa) ---
    \node[draw=gray!35, fill=white, rounded corners=3pt, line width=0.5pt, anchor=south, inner sep=4.5pt] at (4.5, 4.85) {
        \scriptsize
        \tikz[baseline=-0.5ex]{\draw[ndcgline] (0,0) -- (0.45,0) plot[mark=+, mark size=1.8pt] coordinates {(0.225,0)};}~NDCG@5 (asse sinistro) \qquad
        \tikz[baseline=-0.5ex]{\draw[lossline] (0,0) -- (0.45,0) plot[mark=x, mark size=1.6pt] coordinates {(0.225,0)};}~Pairwise Loss (asse destro)
    };

\end{tikzpicture}
\caption{Divergenza empirica tra funzione di perdita pairwise e metrica di ranking NDCG@5 durante l'addestramento di RankNet (Taylor et al.~\cite{taylor2006}).}
\label{fig:ltr-pairwise-loss-vs-ndcg}
\end{figure}

L'andamento evidenzia il limite intrinseco dell'approccio pairwise rispetto all'ottimizzazione posizionale: mentre la perdita continua a decrescere monotonicamente (asse destro in ottanio), la metrica NDCG@5 (asse sinistro in bordeaux) subisce un progressivo degrado dopo aver raggiunto il picco attorno all'epoca 50--70. Ci\`o dimostra che la minimizzazione globale delle coppie errate non garantisce l'ottimizzazione dei primi documenti restituiti all'utente, motivando storicamente l'introduzione di approcci sensibili alla posizione come LambdaRank.

La stabilizzazione piatta di NDCG@5 osservabile tra l'epoca 150 e 200 riflette la natura discreta della metrica a gradini: finch\'e la permutazione relativa dei primi 5 documenti non varia, il valore di NDCG rimane rigorosamente costante anche a fronte di continue variazioni dei punteggi continui prodotti dal modello.

\subsection{Approccio Listwise e l'Intrattabilit\`a dell'Ottimizzazione Diretta}
\label{subsec:approccio-listwise-intrattabilita}

Il paradigma \textbf{Listwise} si propone di risolvere la discrepanza affrontando l'intero elenco dei candidati come un'unica entit\`a indivisibile e ottimizzando direttamente le metriche posizionali ufficiali (NDCG, MAP).

Tuttavia, l'applicazione standard della Discesa del Gradiente si scontra con una barriera matematica:

\begin{theorem}[Non-Derivabilit\`a dell'Operatore di Ranking]
Sia $\mathbf{s} = \langle s_1, s_2, \dots, s_n \rangle \in \mathbb{R}^n$ il vettore dei punteggi continui prodotti dal modello di ranking $h_{\mathbf{\theta}}$. La permutazione o posizione di rango $\pi(i)$ del documento $d_i$ \`e determinata dall'operatore di ordinamento discreto:
\begin{equation}
    \pi(i) = 1 + \sum_{j \neq i} \mathbb{I}(s_j > s_i),
    \label{eq:rank-indicator-sum}
\end{equation}
dove $\mathbb{I}(\cdot)$ indica la funzione indicatrice booleana.  
Poich\'e la funzione indicatrice \`e una funzione a gradino costante a tratti con derivata nulla quasi ovunque e punti di discontinuit\`a non derivabili ogni volta che $s_j = s_i$, ne consegue che:
\begin{equation}
    \frac{\partial \pi(i)}{\partial s_k} = 0 \quad \text{quasi ovunque}, \quad \text{e indifferenziabile nei punti di pareggio.}
\end{equation}
Di conseguenza, qualsiasi metrica $M(\pi, \mathbf{y})$ definita sui ranghi posizionali (come NDCG@$k$ o MAP) possiede gradiente nullo o non definito rispetto ai parametri $\mathbf{\theta}$ del modello:
\begin{equation}
    \nabla_{\mathbf{\theta}} \text{NDCG}@k = \mathbf{0}.
\end{equation}
\end{theorem}

L'impossibilit\`a di applicare il calcolo differenziale diretto alle funzioni di ranking ha aperto la strada a due strategie di risoluzione storiche:
\begin{enumerate}
    \item L'ottimizzazione di funzioni di perdita differenziabili surrogate che fungono da limite superiore per NDCG (ad esempio \textbf{ListNet}, \textbf{ListMLE});
    \item La formulazione di gradienti virtuali corretti dall'impatto sul ranking, alla base della teoria di \textbf{LambdaRank} e dell'algoritmo allo stato dell'arte \textbf{LambdaMART}~\cite{burges2010}, discussi nel Capitolo successivo.
\end{enumerate}

\FloatBarrier
\section{Caso di Studio: BM25F e Ottimizzazione di Ranking Multi-Campo}
\label{sec:bm25f-ottimizzazione-strutturata}

Per comprendere concretamente come i principi di Learning to Rank trovino applicazione nella calibrazione di funzioni di ranking strutturate, questa sezione esamina in dettaglio l'estensione \textbf{BM25F} (Robertson et al., Zaragoza et al.~\cite{taylor2006}).

\subsection{Fondamenti e Motivazione di BM25F}
\label{subsec:fondamenti-bm25f}

Nei documenti complessi (articoli scientifici, brevetti, pagine web), il testo \`e suddiviso in campi o zone semi-strutturate $s \in \mathcal{S}$:
\begin{itemize}
    \item Per un articolo accademico: $\mathcal{S} = \{\text{titolo}, \text{abstract}, \text{corpo}, \text{autori}\}$;
    \item Per una pagina Web: $\mathcal{S} = \{\text{titolo}, \text{url}, \text{body}, \text{anchor text}\}$.
\end{itemize}

L'approccio ingenuo consisterebbe nel calcolare lo score Okapi BM25 standard separatamente su ciascun campo $s$ e combinare linearmente i punteggi parziali ottenuti:
\begin{equation}
    s_{\text{ingenuo}}(q, d) = \sum_{s \in \mathcal{S}} w_s \cdot \text{BM25}_s(q, d).
\end{equation}
Tale combinazione soffre di un difetto intrinseco: applica la saturazione non lineare del termine indipendentemente su ogni singolo campo. Di conseguenza, se un termine compare due volte nel titolo e due volte nel testo principale, esso subisce due saturazioni separate, anzich\'e saturare una volta sola in base all'evidenza complessiva ponderata del termine all'interno del documento.

L'algoritmo \textbf{BM25F} (*BM25 with Fields*) adotta la corretta impostazione algebrica: \textbf{combina linearmente le frequenze ponderate dei campi all'interno del termine prima di applicare la trasformazione di saturazione asintotica}.

\subsection{Formulazione Matematica di BM25F}
\label{subsec:formulazione-matematica-bm25f}

Sia $f_{t, s}$ la frequenza del termine $t$ nel campo $s$ del documento $d$.  
Si definiscono:
\begin{itemize}
    \item $w_s > 0$: il peso d'importanza associato al campo $s$;
    \item $l_s$: la lunghezza effettiva (in numero di termini) del campo $s$ nel documento $d$;
    \item $L_s$: la lunghezza media del campo $s$ sull'intera collezione di documenti;
    \item $b_s \in [0, 1]$: il parametro di normalizzazione della lunghezza specifico per il campo $s$.
\end{itemize}

La lunghezza normalizzata del campo $s$ \`e calcolata come:
\begin{equation}
    B_s = (1 - b_s) + b_s \cdot \frac{l_s}{L_s}.
    \label{eq:field-length-norm}
\end{equation}

La \textbf{pseudo-frequenza effettiva ponderata} del termine $t$ nel documento $d$, denotata con $\tilde{f}_t(d)$, \`e la combinazione convessa:
\begin{equation}
    \tilde{f}_t(d) = \sum_{s \in \mathcal{S}} w_s \cdot \frac{f_{t, s}}{B_s} = \sum_{s \in \mathcal{S}} \frac{w_s \cdot f_{t, s}}{(1 - b_s) + b_s \cdot \frac{l_s}{L_s}}.
    \label{eq:bm25f-weighted-tf}
\end{equation}

Il punteggio globale BM25F della coppia $(q, d)$ applica la saturazione non lineare basata sul parametro $k_1$:
\begin{equation}
    s_{\text{BM25F}}(q, d) = \sum_{t \in q} \text{IDF}(t) \cdot \frac{\tilde{f}_t(d)}{k_1 + \tilde{f}_t(d)},
    \label{eq:bm25f-final-score}
\end{equation}
dove l'inverso della frequenza documentale $\text{IDF}(t)$ \`e calcolato convenzionalmente come:
\begin{equation}
    \text{IDF}(t) = \log \left(\frac{N - df_t + 0.5}{df_t + 0.5}\right).
\end{equation}

\subsection{Esplosione dei Parametri Liberi e Necessit\`a di Learning to Rank}
\label{subsec:esplosione-parametri-bm25f}

Mentre Okapi BM25 possiede soltanto $2$ parametri liberi scalari ($\mathbf{\theta} = \langle k_1, b \rangle$), la funzione BM25F definita su $|\mathcal{S}|$ campi possiede uno spazio di calibrazione di dimensionalit\`a:
\begin{equation}
    |\mathbf{\theta}_{\text{BM25F}}| = 2|\mathcal{S}| + 1.
\end{equation}
Per ciascun campo $s \in \mathcal{S}$, il sistema richiede di fissare:
\begin{itemize}
    \item Il peso d'importanza del campo $w_s$;
    \item Il parametro di normalizzazione della lunghezza $b_s$;
\end{itemize}
a cui si aggiunge il parametro globale di saturazione $k_1$ (oppure parametri $k_{1, s}$ dedicati nelle formulazioni generalizzate).

Se un documento web include ad esempio $S=4$ campi (titolo, body, url, anchor), lo spazio di configurazione comprende $2 \times 4 + 1 = 9$ parametri liberi reali continui.

Una ricerca esaustiva a griglia (\textit{grid search}) risulta proibitiva: campionando appena $10$ valori per ciascuna dimensione, si generano $10^9 = 1$ miliardo di combinazioni, ciascuna richiedente la valutazione di centinaia di query di test su collezioni di scala industriale.  
Risulta dunque mandatorio \textbf{inquadrare BM25F nel paradigma di Learning to Rank}:
\begin{itemize}
    \item \textbf{Spazio delle Ipotesi}: La famiglia di funzioni $h_{\mathbf{\theta}}(q, d) = s_{\text{BM25F}}(q, d; \mathbf{\theta})$, con $\mathbf{\theta} = \langle \{w_s, b_s\}_{s \in \mathcal{S}}, k_1 \rangle$;
    \item \textbf{Dati di Addestramento}: Query annotate con rilevanza graduata $y \in \{0, 1, 2, 3, 4\}$;
    \item \textbf{Metrica di Valutazione}: NDCG@10 (con guadagni esponenziali $2^{y} - 1$ per accentuare l'eccellenza al vertice);
    \item \textbf{Strategia di Ottimizzazione}: Calibrazione automatica dei parametri $\mathbf{\theta}$ mediante discesa di gradiente guidata da una loss surrogata differenziabile (RankNet o gradienti $\lambda$), come formalizzato pionieristicamente da Taylor et al.~\cite{taylor2006}.
\end{itemize}

\FloatBarrier
\section{Quadro Sinottico e Considerazioni Conclusive}
\label{sec:quadro-sinottico-conclusioni}

La Tabella~\ref{tab:ltr-foundations-synoptic} sintetizza i tre paradigmi di Learning to Rank e i rispettivi modelli prototipo analizzati nel capitolo:

\begin{table}[H]
\centering
\small
\resizebox{\linewidth}{!}{%
\begin{tabular}{l p{3.0cm} p{3.2cm} p{3.4cm} p{3.6cm}}
\toprule
\textbf{Paradigma} & \textbf{Unit\`a di Input} & \textbf{Loss Function Prototipo} & \textbf{Ottimizzazione} & \textbf{Criticit\`a Fondamentale} \\
\midrule
\textbf{Pointwise} & Singola coppia $(q, d)$ & Mean Squared Error (MSE), Cross-Entropy & Regressione diretta / SGD & Ignora la posizione del documento nella graduatoria finale. \\
\midrule
\textbf{Pairwise} & Coppia di documenti $(d_i, d_j)$ per query $q$ & RankNet Loss ($\log(1 + e^{-\sigma(s_i - s_j)})$), Hinge Loss & Discesa del gradiente continua (SGD) & Tratta le violazioni in testa e in coda con peso identico; diverge da NDCG. \\
\midrule
\textbf{Listwise} & Intero insieme dei candidati per $q$ & Surrogate Loss differenziabili (ListNet, LambdaMART) & Gradienti virtuali $\lambda$ / Surrogate ranking loss & Non-derivabilit\`a dell'operatore discreto di ordinamento ($\nabla_{\mathbf{\theta}} \text{sort} = \mathbf{0}$). \\
\midrule
\textbf{BM25F (Parametrico)} & Documenti semi-strutturati a $S$ campi & Calibrazione dei $2S+1$ parametri tramite LTR proxy & Gradient Descent su loss di coppia / coordinate search & Evita saturazioni premature combinando le frequenze prima della non linearit\`a. \\
\bottomrule
\end{tabular}%
}
\caption{Quadro sinottico comparativo dei paradigmi di Learning to Rank.}
\label{tab:ltr-foundations-synoptic}
\end{table}

Nel prossimo capitolo l'analisi sar\`a estesa alla formulazione rigorosa di \textbf{LambdaRank} e \textbf{LambdaMART}, illustrando la costruzione algebrica dei gradienti virtuali $\lambda_{ij}$ scalati dal differenziale di metrica $|\Delta \text{NDCG}|$ e l'impiego dei Gradient Boosted Decision Trees (GBDT) nella produzione industriale moderna.
